\chapter{The degree homomorphism}\label{chap:degree}

Each double coset splits into finitely many left cosets, and counting them defines a ring
homomorphism from the Hecke ring to $\mathbb{Z}$ that records the size of each operator.

\begin{definition}\label{def:hecke-coset-deg}
\lean{HeckeRing.HeckeCoset_deg}\leanok
\uses{def:hecke-coset, def:decomp-quot}
Let $(G,H,\Delta)$ be a Hecke pair and let $D=HgH$ be a double coset. Writing
$D=\bigsqcup_i Hg_i$ for its decomposition into finitely many left cosets, the
\emph{degree} of $D$ is
\[
  \deg D \;=\; [\,H : H\cap gHg^{-1}\,],
\]
the number of left cosets $Hg_i$ comprising $D$.
\end{definition}

\begin{lemma}\label{lem:smulOrbit-card}
\lean{HeckeRing.smulOrbit_card}\leanok
\uses{def:hecke-coset-deg}
Fix a double coset $D=HgH$ with left-coset representatives $g_i$, and let $H\beta$ be any
left coset with $\beta\in\Delta$. Then the orbit
\[
  \{\,Hg_i\beta : i\,\}
\]
has cardinality equal to $\deg D$.
\end{lemma}

\begin{proof}\leanok
\uses{def:hecke-coset-deg}
The map sending the $i$-th representative to the left coset $H(\beta g_i g)$ is a
surjection onto the orbit by construction, so it suffices to show it is injective. If two
representatives produced the same coset, then after left-cancelling $\beta$ the
corresponding cosets $Hg_i g$ and $Hg_j g$ would meet, hence coincide; this contradicts
the fact that distinct representatives index distinct left cosets in the decomposition of
$D$. Thus the orbit has exactly as many elements as there are left cosets in $D$, namely
$\deg D$ [Sh, \S3.1].
\end{proof}

\begin{definition}\label{def:deg}
\lean{HeckeRing.deg}\leanok
\uses{def:hecke-coset-deg, def:T-single, thm:hecke-ring}
The \emph{degree homomorphism} is the ring homomorphism
\[
  \deg : \mathbb{T} \longrightarrow \mathbb{Z}
\]
determined by sending each basis element $[D]$ to $\deg D$ and extending
$\mathbb{Z}$-linearly; that is, $\deg\bigl(\sum_D a_D\,[D]\bigr)=\sum_D a_D\,\deg D$. It
sends the identity to $1$, is additive, and is multiplicative
[Sh, \S3.1, Prop.~3.3].
\end{definition}

\begin{lemma}\label{lem:deg-mul}
\lean{HeckeRing.deg_mul}\leanok
\uses{def:deg, lem:smulOrbit-card, def:hecke-multiplicity}
For all $f,g$ in the Hecke ring,
\[
  \deg(f\cdot g)=\deg(f)\,\deg(g).
\]
\end{lemma}

\begin{proof}\leanok
\uses{def:deg, lem:smulOrbit-card, def:hecke-multiplicity}
Since $\deg$ is additive, it suffices to verify multiplicativity on a pair of basis
elements $[D_1]$ and $[D_2]$. Expanding the product in the Hecke ring via the structure
constants gives $[D_1]\cdot[D_2]=\sum_d \mu(D_1,D_2;d)\,[d]$, so by definition
\[
  \deg\bigl([D_1]\cdot[D_2]\bigr)=\sum_d \mu(D_1,D_2;d)\,\deg d.
\]
The right-hand side counts, with multiplicity, the left cosets occurring in the product
double coset. Grouping these by the right action of $D_2$ on the left cosets of $D_1$ and
applying the orbit-size identity, each orbit contributes exactly $\deg D_2$ cosets, and
there are $\deg D_1$ orbits; hence the total is $\deg D_1\,\deg D_2$ [Sh, \S3.1].
\end{proof}
